\footnotesize%
\pstart%
\noindent%
Beide \textit{Schedae de impulsu eccentrico} N.\;\ref{58940} und N.\;\ref{58939} wurden von Leibniz eigenhändig auf den 31.\;Okto\-ber\,/\,10.\;November) 1690 datiert.
Der bereits am gemeinsamen Titel erkennbare enge Zusammenhang beider Aufzeichnungen ist auch dadurch belegt, dass ein Kustos am Schluss von N.\;\ref{58940} %
ausdrücklich auf N.\;\ref{58939} verweist (siehe die Beschreibung des Textzeugen, S.\;\pageref{LH_35_11_14_005-006_ueberlieferung}).
Dies rechtfertigt die Entscheidung, beide Konzepte N.\;\ref{58940} und N.\;\ref{58939} einheitlich zu edieren.
\pend%
\pstart%
Im Hinblick auf die relative Chronologie %
ist bemerkenswert, dass die \textit{Scheda} N.\;\ref{58940} (\refpassage{LH_35_11_14_005-006_005r_cntrgyr-1}{LH_35_11_14_005-006_005r_cntrgyr-2}) eine ähnliche Ansicht über den Mittelpunkt der von einem exzentrischen Stoß herrührenden Drehung vertritt wie der vermutlich in Augsburg\indexO{Augsburg} abgefasste Entwurf N.\;\ref{60300} (\refpassage{LH_37_05_044r_alias_dhe-1}{LH_37_05_044r_alias_dhe-2}):
Das \textit{centrum gyrationis} falle mit dem Schwerpunkt des angestoßenen Körpers zusammen.
Diesbezüglich widerlegt N.\;\ref{60300} ausdrücklich die Ansicht, die in den wohl noch Ende 1689 oder Anfang 1690 in Italien\indexO{Italien} entstandenen Texten N.\;\ref{60142+60297} und N.\;\ref{60299} vetreten worden war (siehe hierüber die Datierungsgründe zu N.\;\ref{60300}, S.\;\pageref{LH_37_05_044_datierung}).
Diese ältere Auffassung des Mittelpunkts der Drehung wird in N.\;\ref{58940} und N.\;\ref{58939} nicht mehr erwähnt.
Dies bekräftigt die Annahme, dass die \textit{Schedae de impulsu eccentrico} nach dem Entwurf N.\;\ref{60300} entstanden sein müssen.
\pend
\selectlanguage{latin}%
\frenchspacing%
\normalsize%
\newpage %
